% Part: computability
% Chapter: computability-theory
% Section: equiv-ce-defs

\documentclass[../../../include/open-logic-section]{subfiles}

\begin{document}

\olfileid[tr]{cmp}{thy}{eqc}

\olsection[Hesaplanabilir Biçimde Numaralandırılabilir Kümelerin
  Tanımları]{Hesaplanabilir Biçimde Numaralandırılabilir Kümelerin Eşdeğer
  Tanımları}

Aşağıdaki teorem, bir kümenin hesaplanabilir biçimde numaralandırılabilir
olmasının ne anlama geldiğine ilişkin birkaç önemli eşdeğer ifade verir.

\begin{thm}
\ollabel{thm:ce-equiv}
$S$ bir doğal sayılar kümesi olsun. Aşağıdakiler eşdeğerdir:
\begin{enumerate}
\item\ollabel{case:ce} $S$ hesaplanabilir biçimde numaralandırılabilir.
\item\ollabel{case:ran-pc} $S$, \emph{kısmi} hesaplanabilir bir fonksiyonun
  görüntü kümesidir.
\item\ollabel{case:ran-prim} $S$ boştur ya da ilkel özyinelemeli bir
  fonksiyonun görüntü kümesidir.
\item\ollabel{case:ce-domain} $S$, kısmi hesaplanabilir bir fonksiyonun
  \emph{tanım kümesi}dir.
\end{enumerate}
\end{thm}

\begin{explain}
İlk üç madde, boş olmayan herhangi bir hesaplanabilir biçimde
numaralandırılabilir kümenin hesaplanabilir, kısmi hesaplanabilir ya da ilkel
özyinelemeli bir fonksiyondan herhangi biriyle numaralandırılabileceğini
söyler. Dördüncü madde ise $S$ hesaplanabilir biçimde
numaralandırılabilirse bazı $e$ indisi için
\[
  S=\Setabs{x}{\cfind{e}(x)\fdefined}
\]
olduğunu söyler. Başka bir deyişle $S$, $\cfind{e}$ hesabının durduğu
girdiler kümesidir. Bu nedenle hesaplanabilir biçimde numaralandırılabilir
kümelere bazen \emph{yarı karar verilebilir} denir: Bir sayı kümedeyse
sonunda ``evet'' yanıtını alırsınız; değilse hiçbir zaman ``hayır'' yanıtını
alamazsınız.
\end{explain}

\begin{proof}
Her ilkel özyinelemeli fonksiyon hesaplanabilir ve her hesaplanabilir
fonksiyon kısmi hesaplanabilir olduğundan \olref{case:ran-prim},
\olref{case:ce} maddesini; \olref{case:ce} ise
\olref{case:ran-pc} maddesini gerektirir. $S$ boşsa, hiçbir yerde tanımlı
olmayan kısmi hesaplanabilir fonksiyonun görüntü kümesi olduğuna dikkat edin.
\olref{case:ran-pc} maddesinin \olref{case:ran-prim} maddesini gerektirdiğini
gösterirsek ilk üç maddenin eşdeğerliğini göstermiş oluruz.

$S$, kısmi hesaplanabilir $\cfind{e}$ fonksiyonunun görüntü kümesi olsun.
$S$ boşsa işimiz biter. Değilse $a$, $S$'nin herhangi bir elemanı olsun.
Kleene normal biçim teoremine göre
\[
  \cfind{e}(x)=U(\umin{s}{T(e,x,s)})
\]
yazabiliriz. Özellikle $\cfind{e}(x)\fdefined$ ve değeri $y$ olur ancak ve
ancak $T(e,x,s)$ ile $U(s)=y$ koşullarını sağlayan bir $s$ varsa. $f(z)$
fonksiyonunu şöyle tanımlayalım:
\[
  f(z)=\begin{cases}
    U((z)_1) & \text{$T(e,(z)_0,(z)_1)$ ise}\\
    a & \text{aksi durumda.}
  \end{cases}
\]
$T$ ile $U$ ilkel özyinelemeli olduğundan $f$ de ilkel özyinelemelidir.
\iftag{TMs}{Turing makineleri cinsinden söylersek $z$, $(z)_1$'in $M_e$
makinesinin $(z)_0$ girdisindeki duran bir hesabı olduğu
$\tuple{(z)_0,(z)_1}$ çiftini kodluyorsa $f$ hesabın çıktısını; aksi durumda
$a$'yı döndürür.}

$S$'nin $f$'nin görüntü kümesi olduğunu, yani her $y$ doğal sayısı için
$y\in S$ ancak ve ancak $y$, $f$'nin görüntü kümesindeyse geçerli olduğunu
göstermeliyiz. İleri yönde $y\in S$ olsun. O zaman $y$, $\cfind{e}$
fonksiyonunun görüntü kümesindedir; dolayısıyla bazı $x$ ile $s$ için
$T(e,x,s)$ ve $U(s)=y$ olur. Böylece $y=f(\tuple{x,s})$ elde edilir.
Tersine $y$, $f$'nin görüntü kümesinde olsun. O zaman ya $y=a$ ya da bazı
$z$ için $T(e,(z)_0,(z)_1)$ ve $U((z)_1)=y$ olur. İkinci durumda
$\cfind{e}((z)_0)\fdefined=y$ olduğundan her iki durumda da $y\in S$ olur.

$\cfind{e}(x)\fdefined=y$ gösterimi, ``$\cfind{e}(x)$ tanımlıdır ve $y$'ye
eşittir'' anlamına gelir. Yalnızca $\cfind{e}(x)=y$ de yazabilirdik; fakat ek
ok, kısmi bir fonksiyonla çalıştığımızı anımsatmak bakımından bazen yararlıdır.

\olref{thm:ce-equiv} ispatını tamamlamak için \olref{case:ce} ile
\olref{case:ce-domain} maddelerinin eşdeğerliğini göstermek yeterlidir. Önce
\olref{case:ce} maddesinin \olref{case:ce-domain} maddesini gerektirdiğini
gösterelim. $S$, hesaplanabilir bir $f$ fonksiyonunun görüntü kümesi olsun:
\[
  S=\Setabs{y}{\text{bazı $x$ için }f(x)=y}.
\]
\[
  g(y)=\umin{x}{(f(x)=y)}
\]
olsun. O zaman $g$ kısmi hesaplanabilirdir ve $g(y)$ ancak ve ancak bazı
$x$ için $f(x)=y$ ise tanımlıdır. Başka bir deyişle $g$'nin tanım kümesi,
$f$'nin görüntü kümesidir. \iftag{TMs}{Turing makineleri cinsinden:
$S$ elemanlarını numaralandıran bir $F$ Turing makinesi verildiğinde,
$F$'nin çıktıları arasında belirli bir elemanı arayan; bulursa duran,
bulamazsa sonsuza kadar arayan ve böylece $S$'ye yarı karar veren $G$ Turing
makinesini alın.}{}

Son olarak \olref{case:ce-domain} maddesinin \olref{case:ce} maddesini
gerektirdiğini göstermek için $S$'nin kısmi hesaplanabilir
$\cfind{e}$ fonksiyonunun tanım kümesi olduğunu varsayalım:
\[
  S=\Setabs{x}{\cfind{e}(x)\fdefined}.
\]
$S$ boşsa işimiz biter; değilse $a$, $S$'nin herhangi bir elemanı olsun.
$f$ fonksiyonunu
\[
  f(z)=\begin{cases}
    (z)_0 & \text{$T(e,(z)_0,(z)_1)$ ise}\\
    a & \text{aksi durumda}
  \end{cases}
\]
ile tanımlayalım. Yukarıdaki gibi bir $x$ sayısı $f$'nin görüntü kümesindedir
ancak ve ancak $\cfind{e}(x)\fdefined$, yani $x\in S$ ise.
\iftag{TMs}{Turing makineleri cinsinden: $S$'ye yarı karar veren bir $M_e$
makinesi verildiğinde, bütün olası Turing makinesi hesapları arasında
dolaşıp duran hesaplara karşılık gelen girdileri döndürerek $S$ elemanlarını
numaralandırın.}{}
\end{proof}

\olref{case:ce-domain} maddesi, \olref{thm:ce-equiv} teoreminde
hesaplanabilir biçimde numaralandırılabilir kümeleri numaralandırmak için
elverişli bir yol sağlar: Her $e$ için $W_e$, $\cfind{e}$ fonksiyonunun tanım
kümesini göstersin, yani
\[
  W_e=\Setabs{x}{\cfind{e}(x)\fdefined}.
\]
O zaman $A$ herhangi bir hesaplanabilir biçimde numaralandırılabilir kümeyse
bazı $e$ için $A=W_e$ olur.

Aşağıdaki sonuç, hesaplanabilir biçimde numaralandırılabilir kümelerin bir
başka karakterizasyonunu verir.

\begin{thm}
\ollabel{thm:exists-char}
Bir $S$ kümesi hesaplanabilir biçimde numaralandırılabilirdir ancak ve ancak
\[
  S=\Setabs{x}{\lexists[y][R(x,y)]}
\]
olacak hesaplanabilir bir $R(x,y)$ bağıntısı vardır.
\end{thm}

\begin{proof}
İleri yönde $S$ hesaplanabilir biçimde numaralandırılabilir olsun. O zaman
bazı $e$ için $S=W_e$ olur. Bu $e$ değeriyle $S$'yi
\[
  S=\Setabs{x}{\lexists[y][T(e,x,y)]}
\]
biçiminde yazabiliriz. Ters yönde
$S=\Setabs{x}{\lexists[y][R(x,y)]}$ olsun ve $f$'yi
\[
  f(x)\simeq\umin{y}{R(x,y)}
\]
ile tanımlayalım. O zaman $f$ kısmi hesaplanabilirdir ve $S$, $f$'nin tanım
kümesidir.
\end{proof}

\end{document}
